% Parte de ejemplo «introduction» (clase section): la edita solo su dueño, declarado en roles.yaml.
% Escribe una oración por línea y separa los párrafos con una línea en blanco (02-convencion §4.3).
% La cita de ejemplo apunta a una entrada de ejemplo de bib/references.bib: bórrala al reemplazar ese archivo.
% El título de la sección y su etiqueta sec:introduction están en tex/outline/introduction.tex, del editor del artículo:
% aquí no van \section ni \chapter; usa \subsection y niveles menores (02-convencion REQ-CONV-94).

Este es el texto de la parte de ejemplo del proyecto.
Cada parte vive en su propia carpeta y tiene un solo dueño.
Las figuras van en \texttt{figures} y las tablas en \texttt{tables}, dentro de la carpeta de la parte.

Las referencias se citan con su clave del grupo de Zotero del proyecto \parencite{knuthTeXbook1984}.
